\section{Variable-Delay Real-Time RL}
\label{sec:envs}
% =============================================================================

This section formalises the variable-delay setting and then describes
the two environment families we study.

\subsection{From fixed to variable action delay}

\citet{ramstedt2019rtrl}'s RTMDP fixes the action delay at exactly one
timestep.
The augmented state $\mathbf{x}_t = (s_t, a_t)$ carries the action
selected at $t{-}1$, and the transition kernel is
\begin{equation}
  p(\mathbf{x}_{t+1} \mid \mathbf{x}_t, \hat{a}_t)
  = p(s_{t+1} \mid s_t, a_t)\,\delta(a_{t+1} - \hat{a}_t),
\end{equation}
which passes the newly emitted action $\hat{a}_t$ forward to land at
frame $t{+}1$.
In RTMDP, delay is a property of the system rather than a choice made by
the agent.

We generalise RTMDP by allowing the action delay to vary per decision.
At each meta-step the gating policy $\pigate$ chooses a budget
$K_t \in \Kset$, and the selected action lands $K_t$ frames later.
Because the environment evolves during those $K_t$ frames, we must also
specify what happens during the $K_t{-}1$ intermediate frames. We call
these filler actions. In committed-action environments they are reflex
actions that move the agent; in clock environments they are no-ops that
consume the player's clock.

Each meta-decision can be viewed as selecting an
\emph{option}~\citep{sutton1999between,bacon2017optioncritic} of length
$K_t$: $K_t{-}1$ filler steps drawn from a fixed reflex distribution,
followed by one action selected by MCTS.
The induced process is a discrete-time semi-Markov decision process
(SMDP) with holding times chosen by the meta-policy, with meta-reward
\begin{equation}
  R_t \;=\; \sum_{k=0}^{K_t - 1} \gamma^k \, r_{t+k}
\end{equation}
and per-meta-step discount $\gamma^{K_t}$, as in the standard options
formulation. The Bellman equation is
\begin{equation}
  V(s_t)
  \;=\;
  \mathbb{E}_{K_t \sim \pigate}\!\left[
    R_t + \gamma^{K_t}\, V(s_{t + K_t})
  \right],
  \label{eq:smdp}
\end{equation}
and we adapt Generalized Advantage Estimation~\citep{schulman2015high}
to carry the per-step discount $\gamma^{K_t}$ through the advantage
computation.
In simulation this is a discrete SMDP with deterministic holding time
set by $K_t$; in deployment, MCTS runtime is slightly stochastic, so
the hardware system is closer to a continuous-time SMDP. Appendix
~\ref{app:rtrl_smdp} discusses this distinction.

Across both environment families, planning is cost-aware by
construction: we do not add a separate reward penalty for thinking.
Either the search tree simulates the effect of delay directly, or the
network is trained on observations that include the remaining time
budget.

\subsection{Committed-action environments: Pac-Man, real-time Tetris, Snake}

In committed-action environments, acting and planning happen
concurrently, and cost-awareness is built directly into the MCTS
search.
When the gating policy selects $K$, it reserves $K$ environment frames
for the planner.
The agent emits $K{-}1$ filler actions---reflex argmax actions drawn
from the planner's own policy network, computed via a single forward
pass in roughly two milliseconds---while the full MCTS runs in the
background.
On the $K$-th frame the planner's chosen action is applied.
The MCTS tree rolls out the same $K{-}1$ reflex steps that the agent
will actually execute, so search is performed over the future state in
which the planned action will land.
The cost of deeper planning is therefore exactly $K{-}1$ additional
frames of world progression before the carefully chosen action lands.

We instantiate this realisation in three environments derived from
Jumanji~\citep{bonnet2024jumanji}: Pac-Man, Snake, and a real-time
variant of Tetris.
The source of progression differs across environments---ghost motion in
Pac-Man, gravity in real-time Tetris, body elongation in Snake---but in
all cases a deeper plan arrives later, after the world has changed.
We modify Jumanji Tetris so that gravity advances the piece in real
time; environment-specific details are in Appendix~\ref{app:envs_tetris}.

\paragraph{Equal compute per frame.}
All budget options in $\Kset$ are calibrated to carry the same number
of MCTS simulations per game frame: at budget $K$ the planner runs
$K \times 32$ simulations over $K$ environment frames, so each frame
receives 32 simulations on average regardless of $K$
(Table~\ref{tab:budget}).
Differences in performance therefore reflect \emph{when} compute is
allocated, not the average compute spent per frame.

\begin{table}[h]
  \centering
  \caption{Committed-action environments: all budget options share 32
    MCTS simulations per game frame.}
  \label{tab:budget}
  \small
  \begin{tabular}{lcccc}
    \toprule
    & $K{=}1$ & $K{=}2$ & $K{=}3$ & $K{=}4$ \\
    \midrule
    MCTS simulations       & 32 & 64 & 96 & 128 \\
    Filler actions         & 0  & 1  & 2  & 3   \\
    Simulations per frame  & 32 & 32 & 32 & 32  \\
    \bottomrule
  \end{tabular}
\end{table}

\subsection{Clock environments: Speed Hex and Speed Go}

In clock environments the board does not evolve while a player
deliberates, but each player has an explicit clock that depletes with
thinking time. Filler actions are no-ops, and only the planner's chosen
action affects the board. Cost-awareness is built into the value
network rather than the search tree: both players' remaining time is
part of the observation during self-play, so the value estimates already
price in the cost of thinking. MCTS can therefore simulate only
board-state dynamics while remaining time-aware through the priors and
values it consumes.
We instantiate this realisation in Speed Hex (11$\times$11) and Speed
Go (9$\times$9), both built on \texttt{pgx}; in both cases the central
decision is not only \emph{which} move to play but \emph{how much} of
the remaining clock to spend before playing it.

\paragraph{Simulation-option calibration.}
For clock environments the equal-compute-per-frame constraint does not
apply, so we calibrate the simulation options $\Kset$ separately to
ensure each option represents a meaningfully distinct tradeoff between
inference latency and planning quality.
Calibration details and the resulting option sets---$\{2,8,32,128\}$
simulations for Speed Hex and $\{16,32,64,96\}$ for Speed Go---are in
Appendix~\ref{app:sim_calibration}.


% =============================================================================
